\subsubsection{Bases Teóricas}

\paragraph{1. Transformación Digital y Comunicación Institucional}
La transformación digital rediseña procesos, canales y decisiones institucionales, habilitando comunicación continua, ágil y centrada en usuarios con expectativas crecientes de inmediatez.

\paragraph{2. Software como Servicio (SaaS)}
Sofware en la nube accedido por suscripción, con escalabilidad, actualizaciones centralizadas y menor inversión inicial para organizaciones sin infraestructura propia.

\paragraph{3. Sistemas Conversacionales y Agentes de Autorespuesta}
Los agentes conversacionales automatizan consultas frecuentes, clasifican intenciones y responden consistentemente, reduciendo carga operativa y derivando casos complejos a atención humana.

\paragraph{4. Integración Multicanal: WhatsApp y Telegram}
La integración multicanal unifica interacciones de WhatsApp y Telegram, mejora trazabilidad, evita fragmentación informativa y sostiene experiencias coherentes en distintos puntos de contacto.

\paragraph{5. Arquitectura Orientada a Servicios y APIs}
Una arquitectura orientada a servicios desacopla módulos mediante APIs, favorece interoperabilidad, escalabilidad y mantenimiento evolutivo, incorporando funciones sin rediseñar todo el sistema.

\paragraph{6. Calidad de Servicio en la Atención Digital}
La calidad digital se mide con tiempos de respuesta, consistencia, disponibilidad y satisfacción; su optimización fortalece confianza institucional y eficiencia operativa sostenida.

\paragraph{7. Experiencia de Usuario (UX) en Canales Conversacionales}
La UX conversacional exige mensajes claros, pertinentes y breves, minimizando fricción comunicacional y facilitando transiciones fluidas hacia atención humana cuando corresponda.

\paragraph{8. Seguridad, Privacidad y Protección de Datos}
La plataforma debe garantizar confidencialidad, integridad y disponibilidad, aplicando controles de acceso, cifrado, trazabilidad y cumplimiento normativo sobre datos sensibles conversacionales.

\paragraph{9. Procesamiento de Lenguaje Natural (PLN) en Chatbots}
El PLN permite reconocer intenciones y entidades en mensajes, mejorando precisión de respuestas y habilitando automatización contextual en dominios institucionales específicos.

\paragraph{10. Evaluación del Rendimiento del Sistema}
El rendimiento se evalúa con KPIs como tiempo de respuesta, resolución automática, abandono, satisfacción y disponibilidad, orientando decisiones de mejora continua.

En síntesis, estas bases teóricas articulan tecnología, gestión y servicio para sustentar una plataforma SaaS de mensajería institucional eficiente, segura y escalable.